% ============================================================================
% Fichier  : partie1/P01_C01_gouverner_cybersecurite_grc.tex
% Titre    : Gouverner la Cybersécurité : du Contrôle à la GRC
% Auteur   : MINKA MI NGUIDJOÏ Thierry Emmanuel
% Version  : 1.0 -- novembre 2026
% Corrigés : corriges/Corr_C01.tex
% ============================================================================
% Faits du cas (bible des cas) : BANQUE SANAGA, 38 agences, DG M. FOTSO,
% RSSI Mme ABENA rattachée au directeur des risques, audit interne dirigé par
% Mme EKAMBI ; certificat ISO/IEC 27001:2022 sur le seul centre de données ;
% rançongiciel du 14/09/2026. Faits nouveaux introduits ici (reportés dans
% la bible) : le périmètre du SMSI a été proposé par le consultant, accepté
% par le DSI et signé par le DG sans passer par le conseil ; comité de
% sécurité sans procès-verbaux ; évaluation de gouvernance 2025 : 2/4/2.
% Sources : modèle des trois lignes (IIA, 2020) ; ISO/IEC 27001:2022, art. 5
% et 6.1.2 (paraphrasés) ; ISO/IEC 27014:2020 (processus de gouvernance,
% paraphrasés, à vérifier sur le texte) ; NIST CSF 2.0, fonction GOVERN ;
% COBIT 2019 (domaine EDM) ; loi n° 2010/012, art. 26 al. 1 (texte français
% officiel, Ann. D d'ACSI) ; COBAC R-2016/04 (paraphrasé, sans numéro
% d'article).
% ============================================================================

\chapter{Gouverner la Cybersécurité : du Contrôle à la GRC}
\label{chap:grc}
\chaptermark{Gouverner la cybersécurité}

\epigraph{%
    « La sécurité se décide au conseil,\\
    elle s'exécute dans les services,\\
    elle se vérifie ailleurs. »%
}{Principe du cours}

\niveauChapitre{Fondamental}{%
    Maîtrise du Chapitre~\ref{chap:prologue}%
}

\begin{boxobjectifs}
\begin{itemize}
    \item Définir gouvernance, gestion des risques et conformité, et leur
          articulation.
    \item Distinguer la gouvernance de la sécurité de sa gestion
          opérationnelle.
    \item Appliquer le modèle des trois lignes à la sécurité.
    \item Situer les rôles du conseil, de la direction, du comité de
          sécurité, du RSSI, du propriétaire de risque et du délégué à la
          protection des données.
    \item Choisir et justifier le rattachement d'un RSSI.
    \item Construire et vérifier une matrice RACI de la sécurité.
    \item Évaluer un dispositif de gouvernance par $\mu$, $U$ et $H$.
\end{itemize}
\end{boxobjectifs}

\bigskip

% ============================================================================
\section*{Avant de commencer : qui a décidé du périmètre~?}
% ============================================================================

Douala, 6~octobre 2026. Trois semaines après le rançongiciel, le comité
des risques du conseil d'administration de la BANQUE SANAGA\footnote{Comme
dans tout le cours, la BANQUE SANAGA, ses dirigeants, les dates et les
montants sont fictifs~; les normes et les textes cités sont exacts.} se
réunit en séance extraordinaire. Le président du comité ne demande pas
comment l'attaque s'est produite~: la RSSI l'a déjà expliqué. Il pose une
autre question~: « Qui a décidé que les agences resteraient hors du
périmètre du certificat~? »

Personne ne répond tout de suite. On reconstitue~: en 2024, le consultant qui a
accompagné la certification a proposé le centre de données seul~; le
directeur des systèmes d'information l'a accepté, parce que le calendrier
était serré~; le directeur général a signé la déclaration de périmètre, au
milieu d'autres documents. Le conseil d'administration ne l'a jamais vue.
Une décision qui engageait la sécurité de trente-huit agences a été prise
sans que l'organe qui répond des risques de la banque en soit informé.

\medskip

Ce chapitre explique ce qu'est gouverner la cybersécurité, qui en porte la
responsabilité et comment on le vérifie. Il pose le cinquième réflexe du
prologue~: relier norme, risque et audit au lieu de les mener comme trois
exercices séparés.

\begin{boxalert}
\textbf{Violation typique.} Traiter une décision de sécurité structurante
(périmètre, acceptation d'un risque, budget) comme une décision technique,
prise et signée sans que l'instance qui répond des risques en soit saisie.

\smallskip
\textbf{Règle~:} toute décision qui fixe le niveau de risque accepté par
l'organisation remonte à l'instance qui en répond. Le technicien propose,
la direction arbitre, le conseil approuve et surveille.
\end{boxalert}

\separateur

% ============================================================================
\section{Du Contrôle à la GRC}
\label{sec:c01-grc}
% ============================================================================

\subsection{Trois âges de la sécurité}

La sécurité des systèmes d'information n'a pas toujours été une affaire de
gouvernance. On peut distinguer trois âges, qui coexistent encore dans
beaucoup d'organisations.

\begin{tableaucours}{L{2.6cm}L{4.6cm}L{5.0cm}}{%
    \tableauHeader{Âge} &
    \tableauHeader{Question posée} &
    \tableauHeader{Limite}
}
Technique & Quels outils installer (pare-feu, antivirus, sauvegardes)~? &
Personne ne dit quels risques ces outils traitent, ni lesquels ils
laissent \\
Contrôle & Les contrôles prévus sont-ils appliqués~? & On vérifie ce qui
est prévu, pas ce qui manque \\
Gouvernance & Quels risques l'organisation accepte-t-elle, qui en répond,
et comment le sait-on~? & Exige que la direction s'en saisisse
réellement \\
\end{tableaucours}
\vspace{-0.8em}
\begin{center}{\small\itshape Tableau~1.1. Trois âges de la sécurité des
systèmes d'information}\end{center}

La BANQUE SANAGA était à l'âge du contrôle~: ses contrôles du centre de
données étaient documentés et appliqués, et l'audit de certification l'a
confirmé. Elle n'était pas à l'âge de la gouvernance~: personne, au niveau
du conseil, n'avait décidé quels risques la banque acceptait de laisser
hors du périmètre.

\subsection{Gouvernance, risque, conformité}

\begin{boxdef}
\termecle{Gouvernance de la sécurité}\index{gouvernance de la sécurité}~:
ensemble des responsabilités et des pratiques par lesquelles l'organe de
direction fixe les objectifs de sécurité de l'organisation, décide du
niveau de risque acceptable, attribue les rôles et surveille les
résultats~\cite{iso27014_2020}.

\medskip
\termecle{Gestion des risques}\index{gestion des risques}~: activités
coordonnées pour identifier, analyser, évaluer et traiter les risques, puis
en suivre l'évolution. Le Chapitre~\ref{chap:risque} en définit le
vocabulaire.

\medskip
\textbf{Conformité}~: au sens du prologue, satisfaction prouvée d'une
exigence~; appliquée ici à l'ensemble des obligations qui s'imposent à
l'organisation (lois, règlements, contrats, normes auxquelles elle
s'engage) et à ses propres politiques.

\medskip
\termecle{Gouvernance, risque et conformité}\index{GRC} (\sigle{GRC})~:
approche qui pilote la
gouvernance, la gestion des risques et la conformité comme un seul
dispositif, avec un référentiel commun d'objectifs, de risques, de
contrôles et de preuves.
\end{boxdef}

Les trois termes ne désignent pas trois services. Ils désignent trois
questions, qui se posent sur le même objet~:
\begin{itemize}
    \item la \textbf{gouvernance} répond à «~que voulons-nous, et jusqu'où
          acceptons-nous le risque~?~»~;
    \item la \textbf{gestion des risques} répond à «~qu'est-ce qui peut nous
          empêcher d'y parvenir, et que faisons-nous~?~»~;
    \item la \textbf{conformité} répond à «~respectons-nous ce qui nous est
          imposé et ce que nous nous sommes imposé~?~».
\end{itemize}

L'\textbf{audit} n'est pas une quatrième lettre~: il donne une assurance
indépendante sur les trois. C'est pourquoi le cours l'étudie à part
(Chapitre~\ref{chap:audit}).

\begin{figure}[H]
\centering
\begin{tikzpicture}[
    bloc/.style={draw=CouleurPrimaire, thick, rounded corners,
                 fill=CouleurPrimaire!8, minimum width=3.4cm,
                 minimum height=1.3cm, align=center, font=\small},
    audit/.style={draw=CouleurAccent, thick, rounded corners,
                  fill=CouleurAccent!12, minimum width=11.6cm,
                  minimum height=0.9cm, align=center, font=\small},
    fleche/.style={->, thick, CouleurSecondaire}]
    \node[bloc] (G) at (0,0) {\textbf{Gouvernance}\\objectifs, appétence,\\rôles};
    \node[bloc] (R) at (4.1,0) {\textbf{Risque}\\apprécier,\\traiter, suivre};
    \node[bloc] (C) at (8.2,0) {\textbf{Conformité}\\obligations,\\preuves};
    \draw[fleche] (G) -- node[above, font=\scriptsize]{fixe le seuil} (R);
    \draw[fleche] (R) -- node[above, font=\scriptsize]{choisit les mesures} (C);
    \draw[fleche] (C.south) to[bend left=25]
        node[below, font=\scriptsize]{rend compte} (G.south);
    \node[audit] (A) at (4.1,-2.6) {\textbf{Audit}~: assurance indépendante
        sur les trois};
\end{tikzpicture}
\caption{L'articulation de la gouvernance, du risque et de la conformité}
\label{fig:c01-grc}
\end{figure}

\begin{boxinfo}
\textbf{GRC n'est pas un logiciel.} Le marché vend des « outils GRC ». Ils
peuvent aider à tenir les registres, mais un outil ne remplace ni une
décision du conseil sur l'appétence, ni un propriétaire de risque qui
signe. Une organisation qui achète l'outil sans définir les rôles obtient
des registres mieux présentés, et aussi incomplets. Le
Chapitre~\ref{chap:grc-pratique} montre comment faire vivre la GRC au
quotidien.
\end{boxinfo}

% ============================================================================
\section{Gouverner n'est pas Gérer}
\label{sec:c01-gouverner}
% ============================================================================

\subsection{Deux niveaux de décision}

La distinction entre gouvernance et gestion vient de la gouvernance des
entreprises~; les référentiels du numérique l'ont reprise. COBIT sépare
ainsi les objectifs de gouvernance, réunis dans le domaine
\emph{évaluer, diriger, surveiller}, des objectifs de gestion, qui
planifient, construisent, exploitent et contrôlent~\cite{isaca2018cobitintro}.
L'ISO/IEC~27014 applique la même logique à la sécurité de l'information~:
l'organe de direction \emph{évalue} la situation et les risques,
\emph{dirige} en fixant objectifs et politiques, \emph{surveille} les
résultats, et \emph{communique} avec les parties
prenantes~\cite{iso27014_2020}.

\begin{tableaucours}{L{3.0cm}L{4.6cm}L{4.6cm}}{%
    \tableauHeader{Processus} &
    \tableauHeader{Ce que fait l'organe de direction} &
    \tableauHeader{À la BANQUE SANAGA, avant l'incident}
}
Évaluer & Apprécie la situation de sécurité et les risques au regard de la
stratégie & Le conseil ne recevait aucun rapport de sécurité \\
Diriger & Fixe les objectifs, l'appétence, les politiques et les moyens &
La politique existait~; l'appétence au risque cyber n'était pas fixée \\
Surveiller & Vérifie que les objectifs sont atteints, par des indicateurs
et par l'audit & Le seul retour était le certificat \\
Communiquer & Rend compte aux actionnaires, au régulateur, aux clients &
Communication de crise improvisée en septembre 2026 \\
\end{tableaucours}
\vspace{-0.8em}
\begin{center}{\small\itshape Tableau~1.2. Les processus de gouvernance de
la sécurité (d'après l'ISO/IEC~27014)}\end{center}

\begin{boxdef}
\termecle{Organe de direction}\index{organe de direction}~: personne ou
groupe qui répond, en dernier ressort, de la performance et de la
conformité de l'organisation~: le conseil d'administration dans une
société anonyme, le conseil d'administration ou l'autorité de tutelle dans
un établissement public.

\medskip
\termecle{Direction}\index{direction} (ou management)~: personnes qui
dirigent l'organisation au quotidien, sous l'autorité de l'organe de
direction~: le directeur général et ses directeurs.
\end{boxdef}

\subsection{Ce que la norme attend de la direction}

L'ISO/IEC~27001 consacre son article~5 au \emph{leadership}. En substance,
la direction doit démontrer qu'elle s'implique dans le système de
management~: elle établit une politique de sécurité compatible avec la
stratégie, intègre ses exigences dans les processus de l'organisation,
fournit les ressources, et attribue les responsabilités et les autorités
des rôles de sécurité~\cite{iso27001_2022}. L'article~6.1.2 ajoute que
chaque risque identifié doit avoir un \emph{propriétaire}.

La norme n'exige pas que le conseil d'administration approuve le
périmètre. Elle exige que la direction s'implique~; elle ne dit pas
comment l'organe de direction la surveille. C'est là que la gouvernance
d'entreprise prend le relais~: pour une banque, le règlement COBAC
R-2016/04 confie à l'organe délibérant la responsabilité du dispositif de
contrôle interne et de maîtrise des risques, et à l'organe exécutif sa mise
en œuvre~\cite{cobac2016r04}. Le Chapitre~\ref{chap:droit} précisera ces
obligations.

\subsection{Le cadre NIST~: gouverner d'abord}

Dans sa version~2.0, publiée en 2024, le cadre de cybersécurité du NIST a
ajouté une sixième fonction, \emph{Gouverner}, aux cinq fonctions
historiques (identifier, protéger, détecter, répondre,
rétablir)~\cite{nist2024csf}. Elle couvre six catégories~: le contexte de
l'organisation~; la stratégie de gestion des risques~; les rôles, les
responsabilités et les autorités~; la politique~; la supervision~; la
gestion des risques de la chaîne d'approvisionnement. Le cadre la place au
centre des cinq autres~: sans elle, elles n'ont pas d'objectif. Le
Chapitre~\ref{chap:referentiels} étudie ce cadre et ses correspondances
avec l'ISO/IEC~27001.

\begin{boxscenario}{Opération SANAGA. Une décision orpheline}
\textbf{Constatations.}
\begin{itemize}
    \item La déclaration de périmètre du SMSI porte la seule signature du
          directeur général.
    \item Les procès-verbaux du conseil et de son comité des risques de 2024
          et 2025 ne mentionnent ni le périmètre ni le certificat, sauf pour
          annoncer son obtention.
    \item Aucun document ne fixe l'appétence de la banque au risque cyber.
\end{itemize}

\textbf{Analyse} (confiance élevée). La direction a satisfait l'exigence
de la norme~: elle s'est impliquée et a signé. Mais l'organe de direction
n'a ni évalué, ni dirigé, ni surveillé cette décision. On peut discuter,
comme le prologue, de la prise en compte des interfaces et des dépendances
du périmètre~; mais le défaut premier n'est pas de conformité à
l'ISO/IEC~27001, il est de gouvernance. Un audit de certification pouvait
contester le périmètre~; il ne pouvait pas corriger le silence du conseil,
puisqu'il ne porte pas sur son fonctionnement.
\end{boxscenario}

% ============================================================================
\section{Le Modèle des Trois Lignes appliqué à la Sécurité}
\label{sec:c01-trois-lignes}
% ============================================================================

\subsection{Le modèle}

L'Institut des auditeurs internes (\sigle{IIA}) a publié en 2020 le
\termecle{modèle des trois lignes}\index{modèle des trois
lignes}~\cite{iia2020troislignes}. Il remplace l'ancien modèle des
« trois lignes de défense »~: le mot \emph{défense} a disparu, parce que
la gestion des risques sert aussi à saisir des opportunités, et parce que
les lignes doivent coopérer, pas seulement se protéger les unes des autres.
Le modèle distingue des \emph{rôles}, pas des services~: une même
direction peut tenir un rôle de première ligne pour une activité et de
deuxième ligne pour une autre.

\begin{tableaucours}{L{2.7cm}L{4.5cm}L{5.0cm}}{%
    \tableauHeader{Rôle} &
    \tableauHeader{Ce qu'il fait en sécurité} &
    \tableauHeader{À la BANQUE SANAGA}
}
Organe de gouvernance & Fixe l'appétence, approuve la stratégie de
sécurité, surveille & Conseil d'administration, comité des risques du
conseil \\
Première ligne & Exploite les systèmes et applique les mesures de sécurité
au quotidien & DSI (exploitation, réseau, postes), directeurs métier,
chefs d'agence \\
Deuxième ligne & Fixe le cadre, conseille, surveille et remet en cause la
première ligne & RSSI, direction des risques, conformité, DPD \\
Troisième ligne & Donne une assurance indépendante à l'organe de
gouvernance & Direction de l'audit interne (Mme~EKAMBI) \\
Prestataires d'assurance externes & Assurance complémentaire &
CERTIQUAL AFRIQUE, commissaires aux comptes \\
\end{tableaucours}
\vspace{-0.8em}
\begin{center}{\small\itshape Tableau~1.3. Le modèle des trois lignes
appliqué à la sécurité de la BANQUE SANAGA}\end{center}

Le superviseur, la COBAC, n'apparaît pas dans le tableau~: il n'est pas
un prestataire au service de la banque, mais une autorité extérieure qui
contrôle l'ensemble du dispositif, les trois lignes comprises.

\subsection{Où placer le RSSI~?}

Le RSSI est l'acteur le plus difficile à placer, parce qu'on lui confie
souvent des tâches des deux premières lignes à la fois. S'il fixe la
politique, surveille les indicateurs et rend compte à la direction, il
tient un rôle de deuxième ligne. S'il administre lui-même le pare-feu, il
tient un rôle de première ligne. S'il fait les deux, il surveille son
propre travail.

\begin{boxalert}
\textbf{Violation typique.} Le RSSI administre les équipements de
sécurité, puis rédige le rapport qui dit au comité qu'ils fonctionnent.

\smallskip
\textbf{Règle~:} celui qui exploite une mesure n'en atteste pas
l'efficacité. Si l'organisation est trop petite pour séparer les rôles,
elle compense par une revue indépendante (audit interne, auditeur
externe) et le dit dans son rapport.
\end{boxalert}

\subsection{Ce que le modèle ne dit pas}

Le modèle des trois lignes est un modèle de rôles. Il ne dit pas si les
lignes couvrent tout l'univers des risques. À la BANQUE SANAGA, les trois
lignes existaient et fonctionnaient~; mais aucune ne regardait les agences
sous l'angle de la cybersécurité~: la DSI les exploitait sans appréciation
des risques, la RSSI concentrait ses efforts sur le périmètre certifié, et
l'audit interne n'avait pas inscrit la sécurité des agences à son
programme. Trois lignes bien rangées peuvent laisser le même trou. C'est
ce que mesure l'évaluation de la Section~\ref{sec:c01-mesurer}.

% ============================================================================
\section{Les Acteurs et leurs Responsabilités}
\label{sec:c01-acteurs}
% ============================================================================

\subsection{Huit acteurs}

\begin{tableaucours}{L{3.2cm}L{9.0cm}}{%
    \tableauHeader{Acteur} &
    \tableauHeader{Responsabilité principale}
}
Conseil et comité des risques & Fixe l'appétence au risque cyber, approuve
la stratégie et les décisions structurantes (périmètre, risques acceptés
au-delà d'un seuil), reçoit les rapports du RSSI et de l'audit interne \\
Direction générale & Met en œuvre la stratégie, fournit les ressources,
arbitre entre les directions, signe la politique de sécurité \\
Comité de sécurité & Instance de direction qui réunit les directeurs
métier, la DSI, le RSSI et les risques~; suit les projets de sécurité, les
incidents et les indicateurs, et prépare les décisions \\
RSSI & Propose la politique, coordonne l'appréciation des risques,
surveille les mesures, conseille, rend compte \\
DSI & Exploite les systèmes et les mesures techniques, applique la
politique dans les projets et l'exploitation \\
Propriétaire de risque & Directeur qui a l'autorité et les moyens de
traiter un risque~; décide du traitement et accepte le risque résiduel
dans la limite de l'appétence \\
Audit interne & Évalue le dispositif de façon indépendante et rend compte
au conseil \\
DPD & Veille à la conformité des traitements de données personnelles,
conseille, sert de point de contact avec l'autorité de protection \\
\end{tableaucours}
\vspace{-0.8em}
\begin{center}{\small\itshape Tableau~1.4. Les acteurs de la gouvernance
de la sécurité}\end{center}

\begin{boxdef}
\termecle{Propriétaire de risque}\index{propriétaire de risque}~:
personne qui a la responsabilité et l'autorité de gérer un
risque~\cite{iso27000_2018}. Ce n'est pas le RSSI~: c'est en général le
directeur dont l'activité est exposée, parce que c'est lui qui peut
dégager les moyens de traitement et qui supporte les conséquences.

\medskip
\termecle{Appétence au risque}\index{appétence au risque}~: niveau de
risque qu'une organisation est prête à accepter pour atteindre ses
objectifs. Elle est fixée par l'organe de direction~\cite{coso2017erm}~; le
Chapitre~\ref{chap:risque} montre comment l'exprimer.
\end{boxdef}

\subsection{Le propriétaire du risque R-AG-01}

La fiche du prologue décrivait le risque. Il lui manquait les rubriques de
gouvernance~: qui en répond, qui peut l'accepter, qui le surveille.

\begin{boxobjet}{R-AG-01 · Rubriques de gouvernance}
\textbf{Propriétaire du risque~:} le directeur du réseau d'agences. C'est
lui qui dirige les agences, arbitre leur budget et supporte la fermeture
des guichets.

\textbf{Contributeurs~:} DSI (postes, messagerie, partage de fichiers),
RSSI (analyse et mesures), direction de la communication (crise).

\textbf{Instance d'acceptation~:} le comité des risques du conseil. Tant
que le conseil n'a pas fixé l'appétence, aucun directeur ne peut accepter
seul un risque de cette ampleur~; une fois l'appétence fixée, le comité
reste compétent pour tout résiduel qui la dépasse.

\textbf{Surveillance~:} la RSSI suit trois indicateurs (part des comptes
d'agence avec second facteur, délai d'application des correctifs, succès
des tests de restauration) et les présente chaque trimestre au comité de
sécurité.

\textbf{Assurance~:} la direction de l'audit interne inscrit le réseau
d'agences à son programme 2027.
\end{boxobjet}

\subsection{Rattacher le RSSI}

Le rattachement du RSSI n'est pas un détail d'organigramme~: il détermine
ce qu'il peut dire, et à qui. Aucune norme n'impose un rattachement
unique~; chaque option a ses forces et ses risques.

\begin{tableaucours}{L{3.0cm}L{4.6cm}L{4.6cm}}{%
    \tableauHeader{Rattachement} &
    \tableauHeader{Force} &
    \tableauHeader{Risque}
}
DSI & Proximité technique, mise en œuvre rapide & Le RSSI surveille son
propre supérieur~; la sécurité passe après les délais des projets \\
Direction des risques & Intégration au dispositif de maîtrise des risques~;
deuxième ligne naturelle & Distance avec la technique~; risque d'une
sécurité réduite à des registres \\
Direction générale & Autorité, accès direct aux arbitrages & Charge de la
direction générale~; peut isoler le RSSI des autres directions \\
Conformité & Lien avec les obligations légales & Une sécurité réduite à la
conformité, aveugle aux menaces non réglementées \\
\end{tableaucours}
\vspace{-0.8em}
\begin{center}{\small\itshape Tableau~1.5. Les options de rattachement du
RSSI}\end{center}

À la BANQUE SANAGA, la RSSI est rattachée au directeur des risques. C'est
un choix défendable~: il place la sécurité en deuxième ligne, à distance
de la DSI. Mais ce rattachement ne valait rien sans un accès au comité des
risques du conseil~: les rapports de la RSSI s'arrêtaient au directeur des
risques. Le critère décisif n'est pas l'organigramme, c'est le
\emph{droit d'être entendu} par l'instance qui décide.

\begin{boxinfo}
\textbf{Le délégué à la protection des données.} Le DPD et le RSSI
travaillent sur des objets voisins~: les données personnelles sont aussi
des informations à protéger. Leurs missions diffèrent pourtant~: le RSSI
protège l'information de l'organisation contre les menaces~; le DPD veille
aux droits des personnes dont les données sont traitées. Confondre les
deux fonctions crée un conflit d'intérêts, car le RSSI peut avoir à
défendre un traitement que le DPD devrait contester. La place du DPD dans
la loi n°~2024/017 relative à la protection des données à caractère
personnel est étudiée dans l'UE MSN-2S3-JUR01~; le cours n'y revient que
pour situer les rôles.
\end{boxinfo}

% ============================================================================
\section{La Matrice RACI de la Sécurité}
\label{sec:c01-raci}
% ============================================================================

\subsection{Quatre lettres}

\begin{boxdef}
\termecle{Matrice RACI}\index{matrice RACI}~: tableau qui croise des
activités et des acteurs, et indique pour chaque case le rôle de l'acteur~:
\begin{itemize}
    \item \textbf{R} (\emph{Responsible}, réalise)~: exécute l'activité~;
    \item \textbf{A} (\emph{Accountable}, répond)~: rend compte du résultat,
          approuve, et ne peut pas déléguer cette responsabilité~;
    \item \textbf{C} (\emph{Consulted}, consulté)~: donne un avis avant la
          décision~;
    \item \textbf{I} (\emph{Informed}, informé)~: est tenu au courant après.
\end{itemize}
\end{boxdef}

\begin{boxprocedure}
\textbf{Vérifier une matrice RACI.} Le papier de travail qui examine une
matrice contrôle, ligne par ligne~:
\begin{itemize}
    \item \textbf{un seul A} par activité~: deux personnes qui répondent,
          c'est personne~;
    \item \textbf{au moins un R}~: une activité sans exécutant n'est pas
          faite~;
    \item \textbf{un A qui a l'autorité}~: celui qui répond doit pouvoir
          décider et dégager les moyens~;
    \item \textbf{pas de R et de A d'attestation sur la même tête}~: celui
          qui exécute une mesure n'atteste pas son efficacité~;
    \item \textbf{les décisions de risque au bon niveau}~: l'acceptation
          d'un risque au-delà de l'appétence remonte à l'organe de
          gouvernance~;
    \item \textbf{la preuve}~: pour chaque A, un document signé (procès-verbal,
          décision, rapport) montre que le rôle a été tenu.
\end{itemize}
\end{boxprocedure}

\subsection{La matrice de la BANQUE SANAGA, après l'incident}

La RSSI propose au comité de sécurité la matrice du Tableau~1.6. Les
colonnes sont~: CA (conseil d'administration et son comité des risques),
DG, CS (comité de sécurité), RSSI, DSI, PR (propriétaire de risque, ici le
directeur métier concerné), AI (audit interne) et DPD.

\begin{tableaucours}{L{4.6cm}cccccccc}{%
    \tableauHeader{Activité} &
    \tableauHeader{CA} & \tableauHeader{DG} & \tableauHeader{CS} &
    \tableauHeader{RSSI} & \tableauHeader{DSI} & \tableauHeader{PR} &
    \tableauHeader{AI} & \tableauHeader{DPD}
}
Fixer l'appétence au risque cyber & A & R & C & C & -- & C & I & -- \\
Approuver la politique de sécurité & I & A & C & R & C & C & I & C \\
Fixer le périmètre du SMSI & A & R & C & R & C & C & I & I \\
Apprécier les risques d'une activité & -- & I & I & R & C & A & I & C \\
Accepter un risque résiduel sous le seuil & -- & I & I & C & -- & A/R & I & -- \\
Accepter un risque résiduel au-delà du seuil & A & R & C & C & -- & C & I & -- \\
Exploiter les mesures techniques & -- & -- & I & C & A/R & -- & -- & -- \\
Surveiller l'efficacité des mesures & I & I & I & A/R & C & I & -- & -- \\
Notifier un incident majeur au régulateur & I & A & I & R & C & C & I & C \\
Auditer le dispositif de sécurité & I & I & I & C & C & C & A/R & -- \\
\end{tableaucours}
\vspace{-0.8em}
\begin{center}{\small\itshape Tableau~1.6. Matrice RACI de la sécurité de
la BANQUE SANAGA (proposition de la RSSI)}\end{center}

Deux lignes demandent une explication. Pour « fixer le périmètre du
SMSI », le DG et la RSSI réalisent (R) et le conseil répond (A)~: c'est
précisément la décision qui avait échappé au conseil. Pour « auditer le
dispositif », l'audit interne répond seul~; la RSSI n'est que consultée,
parce qu'elle ne peut pas auditer ce qu'elle surveille. La mention
«~A/R~» signale un acteur qui répond de l'activité et la réalise par ses
propres équipes~: elle est normale pour l'exploitation (DSI) ou l'audit
(audit interne), à condition que l'attestation d'efficacité vienne d'un
autre acteur.

\begin{boxinfo}
\textbf{Une matrice se lit aussi par colonne.} Une colonne sans aucun A
désigne un acteur qui ne répond de rien~; une colonne chargée de A désigne
un goulet d'étranglement. Dans le Tableau~1.6, la colonne CA porte trois A,
tous sur des décisions de niveau de risque~: c'est le signe que le conseil
gouverne. Un conseil sans A en sécurité ne gouverne pas la sécurité.
\end{boxinfo}

% ============================================================================
\section{Mesurer la Gouvernance}
\label{sec:c01-mesurer}
% ============================================================================

\subsection{Huit questions}

Une gouvernance se décrit dans des documents~; elle se vérifie dans leur
usage. Après l'incident, Mme~EKAMBI évalue la gouvernance de la sécurité
telle qu'elle fonctionnait en 2025, par huit questions. Chaque réponse est
cotée \OUI{} (prouvé), \NON{} (absent ou contredit par les pièces) ou
\NEK{} (impossible à établir faute de pièce).

\begin{tableaucours}{cL{6.6cm}cL{3.4cm}}{%
    \tableauHeader{N°} &
    \tableauHeader{Question} &
    \tableauHeader{Cote} &
    \tableauHeader{Pièce}
}
G1 & Une politique de sécurité est-elle approuvée par la direction
générale~? & \OUI{} & Politique signée, 2024 \\
G2 & Le périmètre du SMSI a-t-il été approuvé par le conseil ou son comité
des risques~? & \NON{} & Procès-verbaux 2024--2025 \\
G3 & L'appétence au risque cyber est-elle fixée par le conseil~? & \NON{} &
Aucun document \\
G4 & Chaque risque du registre a-t-il un propriétaire désigné~? & \NON{} &
Registre 2025, sans colonne propriétaire \\
G5 & Le RSSI est-il indépendant de la DSI~? & \OUI{} & Organigramme,
fiche de poste \\
G6 & Le comité de sécurité s'est-il réuni au moins deux fois en 2025~? &
\NEK{} & Aucun procès-verbal retrouvé \\
G7 & Un rapport de sécurité a-t-il été présenté au comité des risques du
conseil en 2025~? & \NEK{} & Ordre du jour sans annexe \\
G8 & L'audit interne a-t-il couvert la sécurité dans son programme
2024--2025~? & \NON{} & Programmes d'audit \\
\end{tableaucours}
\vspace{-0.8em}
\begin{center}{\small\itshape Tableau~1.7. Évaluation de la gouvernance de
la sécurité en 2025}\end{center}

\begin{boxlecture}{Gouvernance de la sécurité de la BANQUE SANAGA en 2025}
\textbf{Constatation.} 2~\OUI{}, 4~\NON{}, 2~\NEK{} sur 8 questions~:
\[
\mu = \tfrac{2}{8} = 0{,}25~;\quad U = \tfrac{2}{8} = 0{,}25~;\quad
H = -\left(2 \times 0{,}25\log_2 0{,}25 + 0{,}5\log_2 0{,}5\right)
  = 0{,}5 + 0{,}5 + 0{,}5 = 1{,}50~\text{bit}.
\]

\textbf{Analyse} (confiance élevée sur le calcul et sur les \NON{},
moyenne sur la lecture des \NEK{}). La gouvernance existait sur le papier
(politique, RSSI indépendant) et manquait dans les décisions~: périmètre,
appétence, propriétaires, assurance. Les deux \NEK{} portent sur la
\emph{surveillance}~: rien ne prouve que l'organe de direction ait été
informé. Un \NEK{} sur la surveillance n'est pas neutre~: un conseil qui
ne peut pas prouver qu'il a été informé ne peut pas prouver qu'il a
gouverné.

\medskip
\hfill\lecture{0{,}25}{0{,}25}{1{,}50}
\end{boxlecture}

\subsection{Ce que la mesure ajoute}

Dans l'évaluation des dix contrôles clés, le centre de données certifié
obtenait une conformité de 0,80 (prologue, Tableau~P.3). La gouvernance
qui l'encadrait obtient 0,25.
Les deux chiffres ne se contredisent pas~: ils ne mesurent pas le même
objet. Le premier dit que les contrôles du périmètre sont appliqués~; le
second dit que personne, au niveau qui décide, n'a choisi ce périmètre en
connaissance de cause. Le cinquième réflexe du cours, relier norme, risque
et audit, consiste à lire ces deux niveaux ensemble.

% ============================================================================
\section{Le Cadre Camerounais}
\label{sec:c01-cadre}
% ============================================================================

Deux textes donnent à la gouvernance de la sécurité une assise juridique
au Cameroun. Le Chapitre~\ref{chap:droit} les étudie en détail.

\begin{boxjuris}
\textbf{\loicam{n°~2010/012}, article~26, alinéa~1}~\cite{loi2010012}~:

\medskip
\textit{« Les exploitants des systèmes d'information prennent toutes les
mesures techniques et administratives afin de garantir la sécurité des
services offerts. A cet effet, ils se dotent de systèmes normalisés leur
permettant d'identifier, d'évaluer, de traiter et de gérer continûment les
risques liés à la sécurité des systèmes d'information dans le cadre des
services offerts directement ou indirectement. »}

\medskip
\textbf{Portée pour la gouvernance.} La loi impose aux exploitants de
systèmes d'information un dispositif \emph{normalisé} et \emph{permanent}
de gestion des risques de sécurité, pour les services fournis
« directement ou indirectement ».

\medskip
\textbf{Analyse} (lecture du cours, confiance moyenne en l'absence de
jurisprudence connue). Le texte ne prévoit pas de restreindre ce
dispositif à une partie des services. Un dispositif qui n'apprécie les
risques que sur le centre de données répond donc mal à cette obligation,
quel que soit le certificat.
\end{boxjuris}

Pour les établissements de crédit, le règlement COBAC R-2016/04 organise
le contrôle interne et la maîtrise des risques, y compris les risques liés
aux systèmes d'information, et en répartit la responsabilité entre
l'organe délibérant et l'organe exécutif~\cite{cobac2016r04}. Le défaut de
gouvernance de la BANQUE SANAGA n'est donc pas seulement une faiblesse de
bonne pratique~: il expose la banque à une observation du superviseur.

% ============================================================================
\section{Ce que la Gouvernance demande à chaque Métier}
\label{sec:c01-metiers}
% ============================================================================

\begin{boxmetier}{RSSI}
\textbf{Obtenir le droit d'être entendu}
\begin{enumerate}
    \item Demander un point permanent à l'ordre du jour du comité des
          risques du conseil, au moins deux fois par an.
    \item Présenter peu d'indicateurs, reliés aux risques et à l'appétence,
          et dire ce qu'on ignore (les \NEK{}) autant que ce qu'on sait.
    \item Se dégager de l'exploitation des outils de sécurité, ou faire
          revoir par un tiers ce qu'on exploite soi-même.
\end{enumerate}
\end{boxmetier}

\begin{boxmetier}{Responsable GRC}
\textbf{Un seul référentiel}
\begin{enumerate}
    \item Tenir à jour la matrice RACI et la faire approuver par le comité
          de sécurité.
    \item Vérifier que chaque risque du registre a un propriétaire, et que
          chaque acceptation au-delà du seuil porte la décision de
          l'instance compétente.
    \item Relier registre des risques, déclaration d'applicabilité et
          programme d'audit, pour qu'aucun pan de l'activité n'échappe aux
          trois (Chapitre~\ref{chap:grc-pratique}).
\end{enumerate}
\end{boxmetier}

\begin{boxmetier}{Auditeur en cybersécurité}
\textbf{Auditer la gouvernance, pas seulement les contrôles}

L'auditeur ne s'arrête pas aux mesures techniques. Il demande qui a décidé
du périmètre, qui a accepté chaque risque résiduel, et quelle pièce le
prouve. Une décision structurante sans instance compétente est une
constatation d'audit, même si chaque contrôle est en place.
\end{boxmetier}

\begin{boxmetier}{Gestionnaire des risques cyber}
\textbf{Faire signer les propriétaires}

Le gestionnaire des risques anime l'appréciation, mais ne possède pas les
risques. Il obtient de chaque directeur concerné qu'il reconnaisse son
risque, choisisse un traitement et signe l'acceptation du résiduel. Un
risque que personne n'a signé reste, en pratique, à la charge du RSSI, qui
n'a pas les moyens de le traiter.
\end{boxmetier}

\separateur

\begin{boxresume}
\textbf{Ce qu'il faut retenir de ce chapitre~:}
\begin{itemize}
    \item \textbf{GRC}~: trois questions sur un même objet~: ce que l'on
          veut et jusqu'où l'on accepte le risque, ce qui peut l'empêcher,
          ce que l'on doit respecter. L'audit donne une assurance
          indépendante sur les trois.
    \item \textbf{Gouverner n'est pas gérer}~: l'organe de direction évalue,
          dirige, surveille et communique~; la direction met en œuvre.
    \item \textbf{Trois lignes}~: l'exploitation, la surveillance et
          l'assurance indépendante sont des rôles distincts~; celui qui
          exploite une mesure n'en atteste pas l'efficacité.
    \item \textbf{Le propriétaire de risque} est le directeur qui a
          l'autorité de le traiter, pas le RSSI.
    \item \textbf{Le rattachement du RSSI} compte moins que son droit
          d'être entendu par l'instance qui décide.
    \item \textbf{RACI}~: un seul A par activité, des décisions de risque au
          bon niveau, une preuve pour chaque A.
    \item \textbf{La gouvernance se mesure}~: à la BANQUE SANAGA, 0,25 de
          conformité et 0,25 d'incertitude en 2025, alors que les
          contrôles du périmètre certifié obtenaient 0,80.
\end{itemize}
\end{boxresume}

\begin{boxinfo}
\textbf{Points de contrôle.} Vous avez compris ce chapitre si vous savez~:
\begin{itemize}
    \item dire, sur un exemple, ce qui relève de la gouvernance, de la
          gestion des risques et de la conformité~;
    \item placer les acteurs de la sécurité d'une organisation dans le
          modèle des trois lignes~;
    \item repérer les défauts d'une matrice RACI~;
    \item expliquer pourquoi un défaut de gouvernance échappe à un audit de
          certification.
\end{itemize}
\end{boxinfo}

% ============================================================================
\section*{Exercices}
% ============================================================================

\begin{boxexo}[title={\ding{45}\quad Exercice~1.1. Gouvernance, risque ou conformité~? \niveaubase}]
Classez chacune des activités suivantes dans la gouvernance, la gestion des
risques ou la conformité, et indiquez l'acteur qui en répond à la BANQUE
SANAGA.
\begin{enumerate}[(a)]
    \item Le conseil décide qu'aucune agence ne doit rester fermée plus de
          24~heures pour une cause informatique.
    \item La RSSI anime un atelier qui identifie les scénarios de
          rançongiciel sur le réseau des agences.
    \item La direction juridique vérifie que la banque a bien transmis à la
          COBAC le rapport demandé après la bascule de mars 2026.
    \item Le directeur du réseau d'agences choisit de cloisonner le partage
          de fichiers plutôt que d'accepter le risque.
    \item Le comité des risques du conseil reçoit le rapport trimestriel
          de sécurité.
\end{enumerate}
\end{boxexo}

\begin{boxexo}[title={\ding{45}\quad Exercice~1.2. Les trois lignes de LiZbiZ \niveaumoyen}]
Chez LiZbiZ, M.~KOUNOU, responsable informatique, détient tous les droits
de l'application intégrée~; il administre aussi le pare-feu et rédige
chaque année le rapport de sécurité présenté au directeur général,
M.~ETEME. Mme~NKENG est contrôleuse interne. Il n'y a pas d'audit interne.
\begin{enumerate}[(a)]
    \item Placez chaque acteur dans le modèle des trois lignes, en
          distinguant les rôles qu'il tient effectivement.
    \item Identifiez les deux conflits de rôles les plus graves et
          expliquez le risque que chacun fait courir.
    \item Proposez une organisation réaliste pour une entreprise de cette
          taille, sans créer de direction de l'audit interne.
\end{enumerate}
\end{boxexo}

\begin{boxexo}[title={\ding{45}\quad Exercice~1.3. Corriger une matrice RACI \niveaumoyen}]
Une première version de la matrice de la BANQUE SANAGA contenait les lignes
suivantes (colonnes~: CA, DG, CS, RSSI, DSI, PR, AI)~:
\begin{itemize}
    \item Fixer le périmètre du SMSI~: --, A, --, R, A, --, --~;
    \item Accepter un risque résiduel au-delà du seuil~: --, I, --, A, --,
          C, --~;
    \item Exploiter les mesures techniques~: --, --, --, R, R, --, --~;
    \item Surveiller l'efficacité des mesures~: --, --, I, A, R, --, --~;
    \item Auditer le dispositif de sécurité~: --, I, --, A, --, --, R.
\end{itemize}
\begin{enumerate}[(a)]
    \item Pour chaque ligne, relevez les règles de vérification violées.
    \item Proposez une version corrigée de chaque ligne.
    \item Quelle ligne, si elle était restée en l'état, aurait reproduit le
          défaut à l'origine de l'incident de septembre~2026~?
\end{enumerate}
\end{boxexo}

\begin{boxexo}[title={\ding{45}\quad Exercice~1.4. Un an après \niveauexpert}]
En octobre 2027, Mme~EKAMBI reprend les huit questions du Tableau~1.7. Le
conseil a approuvé un périmètre étendu aux agences et fixé une appétence~;
le registre a une colonne propriétaire, mais trois risques sur
vingt n'ont pas de propriétaire désigné~; le comité de sécurité s'est
réuni trois fois, procès-verbaux à l'appui~; deux rapports ont été
présentés au comité des risques~; l'audit interne a programmé une mission
sur les agences pour 2028. La politique et le rattachement du RSSI n'ont
pas changé.
\begin{enumerate}[(a)]
    \item Cotez les huit questions et justifiez chaque cote.
    \item Calculez $\mu$, $U$ et $H$, et comparez avec 2025.
    \item La question G4 est-elle bien posée~? Proposez une formulation qui
          permette une cote plus informative, et dites ce qu'elle changerait.
    \item Rédigez, en une demi-page, le mandat du comité de sécurité~:
          composition, fréquence, attributions, ce qu'il décide et ce qu'il
          remonte au conseil.
\end{enumerate}
\end{boxexo}

\bigskip

\begin{boxinfo}
\textbf{Pour aller plus loin.}
\begin{itemize}
    \item Le modèle des trois lignes de l'IIA~\cite{iia2020troislignes}.
    \item L'ISO/IEC~27014:2020~\cite{iso27014_2020}, sur la gouvernance de
          la sécurité de l'information.
    \item Le cadre de cybersécurité du NIST, version~2.0~\cite{nist2024csf},
          et sa fonction \emph{Gouverner}.
    \item L'ISO/IEC~27001:2022~\cite{iso27001_2022}, article~5
          (leadership).
\end{itemize}
\end{boxinfo}

\begin{boxinfo}
\textbf{Transition.} Gouverner, c'est décider du niveau de risque que l'on
accepte. Encore faut-il savoir ce qu'est un risque cyber, comment on le
décrit et comment on le cote, et reconnaître quand une cote n'est qu'une
opinion. C'est l'objet du Chapitre~\ref{chap:risque}.
\end{boxinfo}
